\chapter{Proposta}

Este capítulo apresenta a proposta de melhoria para o repositório \textit{interface-web} do AgroMart,
elaborada a partir dos resultados obtidos na análise estática do repositório derivadas da ferramenta SonarQube.
O objetivo é definir ações para elevar a qualidade interna do código, com foco em prevenção de falhas.

%Qual a proposta? Como se forma? O que ela fala?

%Com base no diagnóstico realizado, foram identificados pontos que impactam diretamente a evolução do sistema, como baixa cobertura de testes e
%ocorrências de \textit{code smells}.
%A proposta, portanto, sugere ações para tratar desses problemas de forma estruturada e priorizada.
%
%As ações sugeridas foram definidas considerando impacto e esforço de implementação, de modo a orientar melhorias progressivas e mensuráveis.
%Ao final, espera-se fortalecer a base de testes automatizados e tornar o código mais estável e sustentável para futuras evoluções,
%além de consolidar boas práticas de desenvolvimento com maior foco em qualidade e segurança.

\section{Problemas identificados}

Após análise via SonarQube o principal problema identificado é a baixa cobertura de testes em código-fonte.
Também é notável a presença de alguns \textit{code smells}, mas esses somam apenas trinta, o que o SonarQube categoriza como nota A,
além disso, também são apontados dois pontos de possível falha de segurança, que após analise mais detalhada, são classificados como falsos positivos.

Dessa forma o foco desta proposta está na cobertura de testes do código, visto que essa se encontra em apenas 0,5\%,
indicando que a maior parte do código todo se encontra sem testes automatizados.

%Listar os principais pontos de melhoria observados no contexto analisado.

\section{Ações propostas}

Descrever as intervenções sugeridas para cada problema identificado, com foco
em correção, prevenção e padronização.

\section{Priorização}

Organizar as ações por ordem de impacto ou urgência, facilitando a execução
da proposta.
